\documentclass{article}
% Shared layout. Keep benchmark prose and result-specific plots outside this folder.
\PassOptionsToPackage{numbers}{natbib}
\usepackage[preprint]{neurips_2026}
\usepackage[T1]{fontenc}
\usepackage[hidelinks]{hyperref}
\usepackage{booktabs,amsmath,microtype,pgfplots}
\pgfplotsset{compat=1.17}
\usetikzlibrary{patterns}
\definecolor{okblue}{HTML}{0072B2}
\definecolor{okorange}{HTML}{D55E00}
\pgfplotsset{every axis/.append style={font=\small, tick label style={font=\small}, width=\linewidth, height=5cm, grid=major}}

\newcommand{\PaperTitle}{Benchmark Payment Delegation: Policy Compliance and Useful Action in a Simulated Ledger}
\newcommand{\PaperOrcid}{0009-0005-9964-6841}
\newcommand{\PaperAbstract}{Delegated payment agents must complete permitted work while respecting spending limits and owner authority. We present a synthetic ledger benchmark with 40 public episodes across 22 scenario templates. Native replay separates task success, ledger goal completion, and policy violations. Scripted reference success is 100.0\% [91.2, 100.0] for the good policy and 0.0\% [0.0, 8.8] for the bad policy. A constant OK response still scores 40.0\% [26.3, 55.4], exposing weak evidence of informed refusal. Removing owner calls reduces success to 90.0\% [76.9, 96.0]. A simulated population sweep varies selection between reference scripts and measures scoring sensitivity. All intervals are 95\%; reference rates use Wilson intervals, and the simulated sweep resamples agent means. These results validate execution and expose scoring limits; they do not measure learned models, payment security in deployment, or human judgment.}
\newcommand{\TaskCount}{40}
\newcommand{\ScenarioCount}{22}
\newcommand{\EpisodeCount}{8640}
\newcommand{\ConditionCount}{108}
\newcommand{\PolicyCount}{8}
\newcommand{\CILevel}{95}
\newcommand{\BootstrapDraws}{2000}
\newcommand{\AgentCount}{20}
\newcommand{\GoodSuccess}{100.0\% [91.2, 100.0]}
\newcommand{\BadSuccess}{0.0\% [0.0, 8.8]}
\newcommand{\ConstantSuccess}{40.0\% [26.3, 55.4]}
\newcommand{\ConstantDecline}{88.9\% [67.2, 96.9]}
\newcommand{\ConstantAction}{0.0\% [0.0, 14.9]}
\newcommand{\BadViolation}{90.0\% [76.9, 96.0]}
\newcommand{\BadGoal}{55.0\% [39.8, 69.3]}
\newcommand{\EmptyGoal}{45.0\% [30.7, 60.2]}
\newcommand{\EmptySuccess}{0.0\% [0.0, 8.8]}
\newcommand{\AblationSuccess}{90.0\% [76.9, 96.0]}
\newcommand{\AblationViolation}{5.0\% [1.4, 16.5]}
\newcommand{\AblationDifference}{10.0 [0.0, 25.0]}
\newcommand{\GoodBadDifference}{100.0 [100.0, 100.0]}
\newcommand{\MidSuccess}{48.2\% [44.5, 52.0]}
\newcommand{\MidViolation}{47.3\% [44.0, 50.4]}
\newcommand{\MidGoal}{77.8\% [75.4, 80.3]}
\newcommand{\ChangedCount}{4}

\pgfplotsset{coverage labels/.style={xtick={0,1,2,3,4},xticklabels={buy,decline,buy with approval,refund,dispute}}}

\title{\PaperTitle}
\author{Pooria Arab\\Independent researcher\\ORCID \texttt{\PaperOrcid}}
\begin{document}
\maketitle
\begin{abstract}\PaperAbstract\end{abstract}
\section{Why this matters}
Delegating a card gives an agent authority to change a financial state. A useful agent must spend when a request permits spending. It must also stop when a request exceeds the owner's authority. These duties can conflict with a narrow completion objective. Buying the requested object can violate a cap. Refusing every purchase can preserve the balance while failing every legitimate request. An evaluation must distinguish both outcomes.\par
Payment delegation also extends beyond purchases. Refunds return funds to a specified instrument. Disputes challenge a particular transaction. These actions can conflict even when each tool accepts its arguments. A refund followed by a dispute can create an invalid sequence. A partial refund can exceed the amount the owner requested. Measuring only the final balance loses the history needed to detect these errors.\par
This paper studies a bounded construct: compliance with explicit delegated-card policy while completing the requested ledger action. The benchmark presents fictional merchants, stated spending rules, and state-changing tools. Its evidence consists of executed tool traces and deterministic replay. It does not estimate financial loss, fraud prevalence, customer satisfaction, or legal compliance. No live payment processor participates in the study.\par
The contribution is a measurement design and an audit of its behavior. We separate native task success, ledger goal completion, and recorded policy violations. We use scripted references to test reachable success and failure. Trivial policies reveal whether the score rewards inaction. An owner-call ablation tests a specific authority boundary. An additional simulated population sweep tests sensitivity between reference behaviors. These experiments support narrow claims about the benchmark, not claims about current language models.\par
A central result limits the interpretation of the score. A constant acknowledgement obtains \ConstantSuccess{} native success without completing any payment action. The benchmark accepts that response on ordinary decline tasks. The result does not establish that an agent understands why it must decline. We retain this finding rather than changing the scorer after observing it. Future evaluation must report action and decline outcomes separately.\par
\section{Related work}
PayBench frames delegated payment authority as a balance between useful autonomy and preserving user intent \citep{paybench}. Its design uses payment traps and harmless lookalikes. We share the concern with authority boundaries, but use a smaller ledger environment with explicit policies. Our corpus is not a balanced collection of matched trap and harmless pairs. A merchant lookalike scenario does not make the entire corpus a matched-pair experiment. We therefore avoid a paired-trap score and make no performance comparison with PayBench.\par
The interaction setting also relates to $\tau$-bench \citep{tau}. That benchmark combines tools, policies, and simulated users, then compares final database state with a goal. Here, owner responses follow a fixed scenario rule rather than a learned user simulator. Replay checks the action history as well as the goal state. This distinction matters because a correct final purchase can follow an earlier violation. A successful goal does not erase the violation.\par
AgentDojo evaluates agents that use tools over untrusted data \citep{agentdojo}. Its attack-and-defense setting motivates separating task utility from security outcomes. Our study does not include an adaptive prompt-injection attacker. The fictional merchant catalog contains confusing identities and policy conflicts, but no evaluated attack population. A clean result here therefore supplies no evidence of prompt-injection resistance. The two settings address related authority problems through different interventions.\par
BetterBench emphasizes benchmark design, documentation, and reproducibility \citep{betterbench}. We apply those principles through a frozen task snapshot, declared controls, generated results, and explicit limits. Statistical reporting also needs a stated source of uncertainty \citep{miller}. Our public tasks form a fixed synthetic corpus. Their confidence intervals rely on hypothetical sampling or simulation assumptions. They do not convert an author-selected corpus into a representative sample of commerce.\par
These comparisons define the gap addressed here. General tool evaluation can miss whether a permitted action also satisfies a specific owner's request. Payment-oriented tasks can reward safe inactivity if utility is omitted. A ledger audit can miss attempted violations if it examines only successful charges. Our design makes these distinctions observable, although its communication check remains weak. The experiments establish that weakness alongside the intended separations.\par
\section{Experiment design}
\subsection{Unit of evaluation and information boundary}
Each episode starts from a policy, a merchant catalog, prior transactions, a balance, and monthly spending. The agent receives the policy and a natural-language request. It can inspect balance and merchant identity before acting. Payment tools accept amounts in integer cents. The final response summarizes the action or decline. The grader then replays the recorded calls from the initial state, instead of trusting the agent's account of events.\par
The agent-visible policy specifies purchase and monthly caps, allowed merchant identifiers, blocked categories, the refund window, and recurring-charge permission. A merchant search can resolve a confusing name. A balance read supplies current spending information. The goal and grading criteria remain evaluation data. Generic controls use only visible messages and tool definitions. Scripted references intentionally use the benchmark's answer-bearing metadata. Their results are reference checks, not independent demonstrations of problem solving.\par
The action space contains balance inspection, merchant search, purchase, refund, dispute, and owner consultation. Owner consultation records that the owner was asked. The response approves a purchase only in the designated approval scenario. This is a simplified interaction model. It does not authenticate a human, check the wording of the question, or model delayed consent. Those omissions constrain any interpretation of approval performance.\par
\subsection{Controls and execution}
The study plan precedes the result files in version history. It specifies the frozen benchmark version, hypotheses, policies, stopping rule, and exclusions. This is local preregistration, not an independently timestamped registry record. We neither deposited a plan with a venue nor ran a human study. Source inspection informed the plan; no study result informed the registered conditions.\par
We execute \PolicyCount{} scripted conditions. Empty returns no final content. Constant returns ``OK.'' Echo repeats the request. Refuse returns a generic refusal. First-tool calls the first advertised tool once and then returns ``OK.'' Packaged-good and packaged-bad execute the supplied reference scripts. No-owner removes owner calls from packaged-good, preserves other calls and text, and drops empty tool-only turns. That last condition changes consultation behavior without relaxing ledger enforcement.\par
All conditions use the same public runner and native grader. The turn and tool budgets are fixed before collection. Infrastructure failures halt the study rather than becoming scored failures. Empty responses, refusals, and unsuccessful actions remain in the denominator. Each condition runs twice. We compare semantic traces after excluding timing fields. Repeated executions test reproducibility; they do not increase the number of independent task observations.\par
\subsection{Additional simulated population sweep}
The simulated experiment defines a population of artificial script selectors. For each simulated agent and task, a hash of the agent seed and task identifier produces a fixed uniform draw. A selection parameter $p$ chooses the good reference script when the draw is below $p$. Otherwise, the simulated agent executes the bad reference script. Selection occurs for the whole episode, so the simulated agent does not switch policies midway through a transaction sequence.\par
The simulated sweep uses the fixed parameter values shown in Table~\ref{tab:simulation}, with \AgentCount{} simulated agents per value. It uses the same draws across parameter values. This coupling makes simulated differences paired and prevents unrelated random draws from obscuring the intervention. It also makes monotonic selection an intentional property of the design. The experiment tests score sensitivity to known behavior mixtures, not emergent capability or learning.\par
The simulated selector reads reference scripts and cannot serve as a deployable policy. Its parameter has no measured correspondence to intelligence, model size, or reasoning ability. A result near $p$ is expected when the good reference always passes and the bad reference always fails. The value of this experiment is narrower: it verifies that executed episodes reproduce the predicted mixture and that violations remain distinct from goal completion.\par
\section{Coverage}
The public corpus contains \TaskCount{} episodes from \ScenarioCount{} scenario templates. Figure~\ref{fig:coverage} reports the distribution by required goal. These counts describe the evaluated corpus, not the prevalence of payment requests in the world. Several templates appear as related variants, so episodes do not constitute independent semantic designs. We retain their relationship when constructing scenario-cluster intervals.\par
\begin{figure}[ht]\centering
\begin{tikzpicture}\begin{axis}[coverage labels,ybar,ymin=0,ylabel={Public episodes},xlabel={Required goal},nodes near coords,bar width=22pt]
\addplot[fill=okblue] table[x=x,y=n,col sep=comma]{data/coverage.generated.csv};
\end{axis}\end{tikzpicture}
\caption{Coverage is uneven across goal types. These are fixed corpus counts, not estimated population rates. A decline-heavy corpus can reward inaction unless useful action is reported separately.}\label{fig:coverage}\end{figure}

Purchase cases cover approved merchants, confusing merchant names, spending thresholds, recurring charges, available balance, and mixed requests. The mixed request requires the permitted part while withholding the blocked part. This provides a useful contrast with a universal refusal policy. A benign purchase at a confusing merchant also prevents the word ``lookalike'' from becoming a synonym for automatic rejection.\par
Decline cases cover over-limit requests, unapproved merchants, blocked categories, split-charge suggestions, exhausted monthly headroom, and prohibited recurring charges. Some require owner consultation even when no purchase should occur. A no-purchase state alone is insufficient for those cases. Other decline cases accept a nonempty response that does not claim a completed purchase, which creates the communication limit exposed by the controls.\par
Refund cases distinguish transaction identity, elapsed time, requested amount, and destination instrument. A valid partial refund leaves an unrefunded remainder. An over-refund can violate the goal even when it remains below the original charge. Dispute cases require the named transaction and a nonempty reason. The reason is not independently judged for truth or legal sufficiency. Thus, dispute success measures a constrained ledger action rather than a valid consumer claim.\par
The corpus does not sample taxes, foreign exchange, chargeback deadlines across jurisdictions, settlement delays, or processor outages. It does not contain a measured distribution of real merchant behavior. Public tasks and scripts can be inspected by anyone evaluating the benchmark. This supports auditing but weakens claims about unseen competence. A separate hidden-world mechanism exists, but this study does not evaluate a private held-out world.\par
\section{Taxonomy of failure types}
\paragraph{Authority failure.} An agent can attempt a purchase beyond a cap, outside the merchant allowlist, or before a required owner consultation. The payment tool blocks specified policy breaches and records the attempt as a violation. The score therefore assesses the proposed action even when no funds move. It does not interpret a processor block as evidence that the agent respected the rule.\par
\paragraph{Identity failure.} Similar merchant names can lead to an unapproved merchant identifier. The agent must use the permitted identity rather than the first plausible name. A search response is evidence about the catalog, not a grant of authority. The allowlist operates on identifiers. This design separates finding a merchant from being permitted to pay that merchant.\par
\paragraph{Goal failure.} A policy-permitted purchase can still be unrequested. The ledger records a violation for the wrong merchant, wrong amount, duplicate goal purchase, or a purchase during a non-purchase task. Some such actions move funds before the environment records the violation. This is an evaluation safeguard, not a claim that every unsafe effect is prevented at execution time. Goal-specific checks rely on evaluator knowledge.\par
\paragraph{Sequence failure.} Refunds and disputes can conflict on the same transaction. Reversing the action order does not make the conflict valid. Refund amounts also accumulate, so a later refund can exceed the requested total. Goal state is evaluated from the current ledger rather than latched when first reached. A later write can invalidate an earlier success.\par
\paragraph{Communication failure.} Silence on a decline task fails unless an owner call supplies the required communication. A first-person completed-purchase claim can also fail when no purchase occurred. Yet generic nonempty content can pass ordinary declines. The check distinguishes silence and selected false claims, not informed refusal. It uses language patterns and cannot establish full semantic agreement with the owner's request.\par
\paragraph{Inaction failure.} A policy that never changes the ledger avoids many violations. It still fails permitted purchases, refunds, and disputes. This is why policy-violation rates need an accompanying utility measure. A low violation rate is compatible with both a capable agent and an inert agent. The action stratum exposes that ambiguity without requiring a new scalar score.\par
\section{Metrics and scoring}
\subsection{Native score and diagnostic measures}
For episode $i$, let $V_i$ denote recorded policy violations and $G_i$ the final ledger goal indicator. Let $A_i$ denote required owner consultation, $N_i$ the no-purchase requirement, $D_i$ dispute correctness when required, and $C_i$ decline communication. Each symbol denotes satisfaction of its applicable requirement, except $V_i$. The native binary score is\par
\begin{equation} S_i=\mathbf{1}\{V_i\leq V_i^{\max}\}\,G_i^{*}A_iN_iD_iC_i,\end{equation}
where $G_i^{*}$ equals $G_i$ when the scenario requires the goal and equals unity otherwise. Requirements that do not apply contribute unity. The public criteria define the permitted violation bound. The study does not relax those criteria. Mean native success is the fraction of episodes with $S_i=1$.\par
We also report the proportion with any recorded violation and the proportion with the ledger goal met. These diagnostics are not substitutes for native success. A blocked purchase may leave a decline goal satisfied while still recording a violation. A valid goal purchase followed by an unrequested action can meet the goal but fail the policy check. The difference between these measures identifies what a single completion score would conceal.\par
Decline and action strata use the task's required goal, fixed before execution. Action includes purchases, approval purchases, refunds, and disputes. We report native success in each stratum without reweighting the corpus. A future balanced score might be useful, but introducing one after observing these results would change the registered analysis. The present paper instead shows the full component table.\par
\subsection{Intervals and dependence}
For a task proportion $\hat p=k/n$, the Wilson interval uses denominator $d=1+z^2/n$, center $(\hat p+z^2/(2n))/d$, and half-width $z\sqrt{\hat p(1-\hat p)/n+z^2/(4n^2)}/d$. We use the normal quantile for \CILevel{}\% coverage. These intervals remain nonzero at observed extremes. They describe hypothetical task sampling and should not be read as measurement noise in a deterministic replay.\par
Task variants share scenario templates. We therefore also resample complete scenario clusters for baseline success. Each sampled cluster contributes all its episodes, and the statistic remains an episode-weighted mean. Paired contrasts resample the same clusters for both policies. This preserves related variants and policy pairing. It does not address selection bias in the choice of scenario templates.\par
The simulated sweep uses a different uncertainty unit. We compute each simulated agent's mean over the fixed corpus, then bootstrap those agent means. Paired simulated contrasts retain the seed correspondence across parameter values. We use \BootstrapDraws{} draws for each bootstrap interval. These simulated intervals describe variation among artificial selectors on this corpus. They cannot estimate variation among learned models or new payment settings.\par
At deterministic endpoints, the bootstrap distribution can collapse. A collapsed interval only says that resampling the observed units yields the same statistic. It does not certify a true population boundary. We retain Wilson intervals beside reference proportions and state endpoint collapse explicitly. No multiplicity-adjusted significance claim or learned-model ranking is made.\par
\section{Results}
\subsection{Reference separation and trivial credit}
The good reference achieves \GoodSuccess{} native success, while the bad reference achieves \BadSuccess{}. Their paired scenario-bootstrap difference is \GoodBadDifference{} percentage points. The collapsed interval follows from uniform task-level separation. This result confirms that the released good and bad scripts reach different scored outcomes. It does not show that independently designed adversarial policies cannot exploit the grader.\par
\begin{table}[ht]\centering\small\begin{tabular}{lll}\toprule
Policy & Success (Wilson) & Violation (Wilson) \\
\midrule
empty & 0.0\% [0.0, 8.8] & 0.0\% [0.0, 8.8] \\
constant & 40.0\% [26.3, 55.4] & 0.0\% [0.0, 8.8] \\
echo & 40.0\% [26.3, 55.4] & 0.0\% [0.0, 8.8] \\
refuse & 40.0\% [26.3, 55.4] & 0.0\% [0.0, 8.8] \\
first-tool & 40.0\% [26.3, 55.4] & 0.0\% [0.0, 8.8] \\
packaged-good & 100.0\% [91.2, 100.0] & 0.0\% [0.0, 8.8] \\
packaged-bad & 0.0\% [0.0, 8.8] & 90.0\% [76.9, 96.0] \\
no-owner & 90.0\% [76.9, 96.0] & 5.0\% [1.4, 16.5] \\
\bottomrule\end{tabular}
\caption{Native success and any policy violation for scripted controls. Values are percentages with task-level Wilson intervals. The policies are not learned models.}\label{tab:baseline}\end{table}
The constant, echo, refuse, and first-tool controls tie at \ConstantSuccess{} overall success. The constant control achieves \ConstantDecline{} on decline tasks and \ConstantAction{} on action tasks. These outcomes show that useful action and decline credit measure different things. A refusal message can satisfy an ordinary decline without proving knowledge of the applicable rule. An acknowledgement and a copied request receive the same credit.\par
The empty control achieves \EmptySuccess{} native success even though its ledger goal rate is \EmptyGoal{}. This is a useful check on communication requirements. A clean no-purchase ledger can satisfy the decline goal, yet silence still fails the complete scoring rule. The result also shows why reporting only ledger goal completion would overstate this control's behavior.\par
\begin{figure}[ht]\centering
\begin{tikzpicture}\begin{axis}[ymin=0,ymax=1.05,ylabel={Episode proportion},xtick={0,1,2,3,4,5,6,7},xticklabels={empty,constant,echo,refuse,first-tool,good,bad,no-owner},x tick label style={rotate=30,anchor=east},legend style={at={(0.5,1.22)},anchor=south,legend columns=2}]
\addplot[okblue,mark=*,error bars/.cd,y dir=both,y explicit] table[x=x,y=success,y error minus=slow,y error plus=shigh,col sep=comma]{data/baseline.generated.csv};\addlegendentry{Success}
\addplot[okorange,mark=square*,error bars/.cd,y dir=both,y explicit] table[x=x,y=violation,y error minus=vlow,y error plus=vhigh,col sep=comma]{data/baseline.generated.csv};\addlegendentry{Any violation}
\end{axis}\end{tikzpicture}
\caption{Scripted controls separate successful action from policy violations. Constant and refusal controls earn the same partial credit. Bars show task-level Wilson intervals; scenario-cluster estimates appear in Table~\ref{tab:cluster}. Lines only connect display positions.}\label{fig:baseline}\end{figure}

The bad reference records violations on \BadViolation{} of episodes while meeting ledger goals on \BadGoal{}. Goal achievement therefore does not imply compliant completion. The bad scripts include blocked attempts and invalid sequences, so the violation measure is broader than successful unauthorized spending. We do not translate those violations into dollars lost or prevented. Their consequences depend on the specific action path.\par
\begin{table}[ht]\centering\small\begin{tabular}{lll}\toprule
Policy & Decline success & Action success \\
\midrule
empty & 0.0\% [0.0, 17.6] & 0.0\% [0.0, 14.9] \\
constant & 88.9\% [67.2, 96.9] & 0.0\% [0.0, 14.9] \\
echo & 88.9\% [67.2, 96.9] & 0.0\% [0.0, 14.9] \\
refuse & 88.9\% [67.2, 96.9] & 0.0\% [0.0, 14.9] \\
first-tool & 88.9\% [67.2, 96.9] & 0.0\% [0.0, 14.9] \\
packaged-good & 100.0\% [82.4, 100.0] & 100.0\% [85.1, 100.0] \\
packaged-bad & 0.0\% [0.0, 17.6] & 0.0\% [0.0, 14.9] \\
no-owner & 88.9\% [67.2, 96.9] & 90.9\% [72.2, 97.5] \\
\bottomrule\end{tabular}
\caption{Native success by required goal stratum, with Wilson intervals. Trivial nonempty policies receive decline credit but complete no action tasks.}\label{tab:strata}\end{table}
\subsection{Owner-call ablation}
Removing owner consultation reduces native success to \AblationSuccess{}. The paired decrease from the good reference is \AblationDifference{} percentage points. The scenario-bootstrap interval includes zero because required consultation covers few scenario clusters. This limits a broad statistical claim even though the individual changed traces have a clear cause.\par
The ablation changes \ChangedCount{} task outcomes. They belong to owner-approved purchase and owner-declined request scenarios and their variants. On an approval purchase, removing the consultation makes the attempted charge violate the owner rule. On an owner-declined request, the absence of consultation fails the explicit requirement even without a purchase. This contrast checks both execution and communication aspects of delegated authority.\par
The ablation's violation rate is \AblationViolation{}, lower than its failure rate. Some missing consultations create failure without a prohibited charge attempt. This difference matters when evaluating controls: a failed task need not be a policy-violating transaction. The ablation removes an agent action, not the enforcement layer. It therefore cannot estimate the effect of turning off processor restrictions or changing real consent procedures.\par
\begin{table}[ht]\centering\small\begin{tabular}{lll}\toprule
Policy & Success (scenario bootstrap) & Ledger goal (Wilson) \\
\midrule
empty & 0.0\% [0.0, 0.0] & 45.0\% [30.7, 60.2] \\
constant & 40.0\% [19.0, 61.5] & 45.0\% [30.7, 60.2] \\
echo & 40.0\% [19.0, 61.5] & 45.0\% [30.7, 60.2] \\
refuse & 40.0\% [19.0, 61.5] & 45.0\% [30.7, 60.2] \\
first-tool & 40.0\% [19.0, 61.5] & 45.0\% [30.7, 60.2] \\
packaged-good & 100.0\% [100.0, 100.0] & 100.0\% [91.2, 100.0] \\
packaged-bad & 0.0\% [0.0, 0.0] & 55.0\% [39.8, 69.3] \\
no-owner & 90.0\% [75.0, 100.0] & 95.0\% [83.5, 98.6] \\
\bottomrule\end{tabular}
\caption{Scenario-cluster intervals for native success and Wilson intervals for ledger goals. Cluster resampling preserves template siblings. Boundary collapse reflects uniform observed outcomes, not certainty outside the corpus.}\label{tab:cluster}\end{table}
\subsection{Simulated population sweep}
The simulated sweep increases successful reference selection while holding tasks and ledger rules fixed. At the midpoint of the simulated sweep, mean success is \MidSuccess{} and mean violation incidence is \MidViolation{}. The same simulated condition meets ledger goals on \MidGoal{}. All displayed simulated intervals resample artificial agent means. This separation reinforces the distinction between finishing a ledger goal and satisfying the complete policy.\par
\begin{figure}[ht]\centering
\begin{tikzpicture}\begin{axis}[ymin=0,ymax=1.05,xmin=0,xmax=1,xlabel={Simulated good-script selection probability},ylabel={Mean episode proportion},legend pos=north west]
\addplot[okblue,mark=*,error bars/.cd,y dir=both,y explicit] table[x=p,y=success,y error minus=slow,y error plus=shigh,col sep=comma]{data/simulation.generated.csv};\addlegendentry{Simulated success}
\addplot[okorange,mark=square*,error bars/.cd,y dir=both,y explicit] table[x=p,y=violation,y error minus=vlow,y error plus=vhigh,col sep=comma]{data/simulation.generated.csv};\addlegendentry{Simulated violations}
\end{axis}\end{tikzpicture}
\caption{The simulated sweep confirms sensitivity to reference-script selection. Bars resample simulated agents, holding public tasks fixed. They do not estimate variation across learned models or real payment environments. Endpoint intervals collapse because every simulated agent selects the same script.}\label{fig:simulation}\end{figure}

\begin{table}[ht]\centering\small\begin{tabular}{lll}\toprule
Simulated mixture $p$ & Success (agent bootstrap) & Violations (agent bootstrap) \\
\midrule
0 & 0.0\% [0.0, 0.0] & 90.0\% [90.0, 90.0] \\
0.25 & 23.6\% [20.6, 26.5] & 69.0\% [66.5, 71.6] \\
0.5 & 48.2\% [44.5, 52.0] & 47.3\% [44.0, 50.4] \\
0.75 & 76.0\% [73.8, 78.3] & 21.6\% [19.7, 23.4] \\
1 & 100.0\% [100.0, 100.0] & 0.0\% [0.0, 0.0] \\
\bottomrule\end{tabular}
\caption{Simulated population sweep. Each row summarizes \AgentCount{} simulated selectors. Intervals bootstrap simulated agent means on the fixed task corpus.}\label{tab:simulation}\end{table}
The simulated success curve follows its constructed mixture. Since every good reference passes and every bad reference fails, success estimates the share of selected good scripts. Its agreement with the selection parameter is a pipeline check, not evidence that payment competence develops smoothly. The simulated violation curve is not a complementary success curve: some failed bad scripts have no recorded violation. That structural difference prevents interpreting every failure as unsafe spending.\par
The simulated contrasts in Table~\ref{tab:contrasts} retain all planned comparisons, including the zero self-contrast. Endpoint intervals collapse because all simulated agents choose the same script. More simulated agents could narrow intervals without adding task diversity. Episode count is not an effective sample size for generalization.\par
\begin{table}[ht]\centering\small\begin{tabular}{ll}\toprule
Simulated mixture $p$ & Success minus simulated $p=0$ \\
\midrule
0 & 0.0 [0.0, 0.0] \\
0.25 & 23.6 [20.6, 26.5] \\
0.5 & 48.2 [44.5, 52.0] \\
0.75 & 76.0 [73.8, 78.3] \\
1 & 100.0 [100.0, 100.0] \\
\bottomrule\end{tabular}
\caption{Paired simulated success differences on a percentage-point scale. The same artificial agent seeds appear in both conditions. These comparisons describe the simulated mixture only.}\label{tab:contrasts}\end{table}
\subsection{Execution evidence and claim boundaries}
Across baseline and simulated conditions, the study executes \EpisodeCount{} episodes. All \ConditionCount{} conditions reproduce their semantic traces on the second run. Native rescoring agrees with recorded scores. These checks support the claim that the reported measurements follow the benchmark's own execution path. They do not provide independent validation of the policy model, because replay uses the same ledger semantics as execution.\par
There are no learned-model scores. Hosted inference was outside the registered local-only plan, and provider availability was not tested. No result should be read as a free-model outage or a comparison of commercial systems. The study incurs no model API charge and uses no real card data. Local execution still consumes compute resources, which are not measured as an energy or total-cost estimate.\par
\section{Roadmap}
The first measurement priority is decline communication. A stronger criterion should require an explicit refusal or an owner consultation tied to the request. It should also accept valid paraphrases and reject irrelevant acknowledgements. Before changing that criterion, a new study should define accepted and rejected cases. A blind annotation exercise could test agreement, but none is claimed here. Published results must remain tied to the present scorer until a versioned change occurs.\par
The next priority is a genuinely paired task design. Each authority violation should have a closely matched permitted request. The same merchant identity, amount context, and request form should differ only at the intended boundary. Such pairs could distinguish indiscriminate caution from rule-sensitive action. The current corpus supports contrasts across goal classes, but it does not justify a causal estimate for any individual trap feature.\par
A further study should compare explicit policy text, agent-side checks, and tool-side enforcement under separate interventions. The present no-owner condition tests only the presence of consultation calls. Removing enforcement would require new environment semantics and an independent audit of possible financial effects. That study should report attempted violations, completed unauthorized effects, and valid utility separately. Combining them into one score would hide different control failures.\par
Learned-model evaluation should freeze prompts, model identifiers, budgets, and refusal handling before collection. Free hosted access would need a spending guard and a review of publication terms. A local model would need its exact weights and inference settings recorded. Either path should retain trivial controls and use a held-out world. A model score without those controls could mistake generic decline credit for delegated-payment competence.\par
Maintenance should preserve previous task and grader versions. A correction must distinguish a task-design change from a scoring change. Re-running old traces under a new grader can measure the latter, but it cannot reconstruct decisions a changed prompt would have produced. Release notes should state which comparisons remain valid. Independent replication is needed before treating this benchmark as a basis for deployment decisions.\par
\section{Caveats}
The construct is narrow and synthetic. Explicit integer-cent policies remove ambiguity that real payment systems often contain. Fictional merchants avoid personal data but do not model the incentives of actual sellers. Owner consultation is deterministic and does not verify human consent. Dispute reasons receive no substantive review. Refund state omits settlement delays and banking operations. These choices support controlled replay while limiting external validity.\par
The reference scripts and grader share design assumptions. Perfect reference success can reveal that intended solutions remain executable, but it cannot prove rubric correctness. The bad reference is one prescribed failure path per task, not an exhaustive attack search. Trivial credit demonstrates one weakness, but no independent expert review establishes the complete set of weaknesses. We therefore describe the study as an execution and measurement audit.\par
Recorded violations combine different events. A blocked purchase, an off-goal completed purchase, and a conflicting refund sequence can each contribute. Their financial severity is not comparable. The paper reports episode incidence rather than adding them into a loss score. Some policy checks use the evaluator's exact goal. Real payment infrastructure may not know that goal, so this scoring model must not be presented as an implemented production control.\par
Statistical uncertainty is conditional on the chosen corpus or simulated population. Scenario-cluster resampling accounts for known siblings but not unobserved families of related requests. Simulated agent resampling captures only the artificial selection mechanism. Narrow intervals can coexist with large external-validity gaps. No confidence interval here measures the probability that a real payment deployment is safe.\par
The text check has its own language boundary. It detects selected first-person completed-purchase claims and clause-level negation. It does not establish whether every final statement is accurate. It also accepts uninformative nonempty responses on ordinary declines. The measured constant-control result makes this limitation concrete. Any downstream use should retain action-stratum reporting and avoid treating decline credit as demonstrated reasoning.\par
The study uses public tasks and no private held-out evaluation. Contamination and deliberate adaptation remain possible. No human baseline, real-user preference study, or production transaction data are available. AI assistance contributed to study scripting and manuscript preparation. That assistance does not constitute independent scientific review. The author remains responsible for the claims and the released evidence.\par
\section{Citation and availability}
Code, frozen task provenance, raw traces, generated data, and reproduction instructions are available at \url{https://github.com/pooriaarab/benchmark-payment-delegation}. Cite this report by title and author, with the version of the evaluation artifact used. The report is a preprint, not a peer-reviewed publication. Paper text and generated figures use CC-BY-4.0. Benchmark code retains MIT, and synthetic task provenance declares CC0-1.0. No new archive DOI is asserted.\par
\section{Authors}
Pooria Arab is an independent researcher. ORCID: \href{https://orcid.org/\PaperOrcid}{\PaperOrcid}. The study evaluates a benchmark maintained by the author. This relationship is a source of potential design bias. The released controls expose a scoring limitation and provide a basis for independent checks, but do not replace outside validation.\par
\small
\bibliographystyle{plainnat}
\bibliography{references}
\end{document}
